\section{A signed interval transfer estimate}\label{lc:transfer}
This section bounds the signed recurrence. It is not a Perron estimate
for its entrywise absolute value. Parametrize the real interval by
\[
 \sigma=-2\cos(2\pi\delta),\qquad u=-4\sin^2(\pi\delta),
                    \qquad0\le\delta\le\tfrac12.
\]
For positive outer sum \(A=x+y\) and \(P=xy\ge0\), define
\begin{equation}\label{lc:continuous-coefficients}
\begin{aligned}
 a_\delta(t)&=4(\cos^2(\pi t)-\sin^2(\pi\delta)),\\
 \rho_\delta(t)&=A\sin(\pi t)+
   \sqrt{(x-y)^2\sin^2(\pi t)+4P\cos^2(\pi\delta)},\\
 \lambda_\delta(t)&=2\sin(\pi t)\rho_\delta(t).
\end{aligned}
\end{equation}
On each compact set of allowed moduli,
\begin{equation}\label{lc:endpoint-size}
 ct\le\rho_\delta(t)\le C,\qquad ct^2\le\lambda_\delta(t)\le C,
                          \qquad0<t\le\tfrac12,
\end{equation}
uniformly in \(\delta\).

\begin{lemma}[Uniform signed transfer]\label{lc:slow-transfer}
For every \(\varepsilon>0\), all sufficiently large \(N\) satisfy
\begin{equation}\label{lc:transfer-estimate}
 |K_{i,j}|\le
 \exp\left\{N\int_{i/N}^{j/N}\log\lambda_\delta(t)\,dt+\varepsilon N\right\},
                      \qquad0\le i\le j\le m.
\end{equation}
The threshold is uniform in endpoints, real spectral parameter,
and each fixed compact set of allowed moduli.
\end{lemma}
\begin{proof}
Normalize \(L_{i,j}=K_{i,j}/\prod_{r=i+1}^js_r\). Its step matrix is
\[
 T_h(t,\delta)=
 \begin{pmatrix}
 2A\sin(\pi t)&4P[\cos^2(\pi(t-h/2))-\sin^2(\pi\delta)]\\1&0
 \end{pmatrix},\quad h=N^{-1},\quad t=jh.
\]
The limiting matrix has real eigenvalues
\[
 A\sin(\pi t)\ \pm\
       \sqrt{(x-y)^2\sin^2(\pi t)+4P\cos^2(\pi\delta)}.
\]
Its spectral radius is \(\rho_\delta(t)\), including coincident roots.
Use the Euclidean operator norm. Fix \(\eta>0\) and restrict to
\(t\in[\eta,1/2]\), where
\(\rho_\delta(t)\ge c_\eta>0\). Cayley--Hamilton gives uniformly
\begin{equation}\label{lc:constant-power}
                    \|T^k\|\le C_\eta(k+1)\rho(T)^k .
\end{equation}
For distinct eigenvalues the power divided difference is a sum of
\(k\) terms, each bounded by \(\rho^{k-1}\). The polynomial identity
extends to coincident eigenvalues by its derivative value. Bounded
entries and the positive lower bound on \(\rho\) give \(C_\eta\).
No eigenbasis or its condition number is used.

Choose a fixed block length \(k\) for which
\(\log(2C_\eta(k+1))/k\) is smaller than a prescribed error.
Within a block, the exact matrices differ from its first limiting
matrix by \(O_k(N^{-1})\). Expanding the finite product one factor
at a time bounds its difference from \(T^k\) by the same order.
For large \(N\) this is at most \(C_\eta(k+1)\rho^k\).
Multiplying block bounds and using uniform continuity of \(\log\rho\)
gives its integral plus the prescribed error. Incomplete blocks have
at most \(2k\) factors and cost \(o(N)\), uniformly in the endpoints.

For \(0<t<\eta\), bounded entries cost at most \(C^{\eta N+O(1)}\).
The omitted logarithmic integral has absolute value at most
\(C\eta\log(2/\eta)\), by \eqref{lc:endpoint-size}.
The exact scalar factors \(s_r\) have the same integrable endpoint
bound, and their logarithmic Riemann sums converge uniformly over
subintervals. Choose \(\eta\) first, then the fixed block length,
and finally \(N\). Restoring these factors proves the estimate.
If a perturbation of the reference matrices is wanted, it must be
chosen after the block length; a sufficiently small fixed perturbation
then fits the same finite-product error. No complex perturbation is
needed for the theorem here.
\end{proof}

\subsection{Unused spectral rows}
The exact diagonals are \(\alpha_r=a_\delta((r+1/2)/N)\).
Uniformly in \(\delta\) and interval endpoints,
\begin{equation}\label{lc:riemann}
 \sum_{r=l}^{k-1}\log|\alpha_r|
 \le N\int_{l/N}^{k/N}\log|a_\delta(t)|\,dt+o(N).
\end{equation}
A sampled zero makes the left side \(-\infty\), which causes no
upper-bound problem. This convention describes a zero product;
no finite division by it is made.

For a proof, truncate the logarithm from below at \(-M\).
The truncated function is jointly continuous, so its Riemann sums
converge uniformly. The factorization
\[
 a_\delta(t)=4\cos(\pi(t+\delta))\cos(\pi(t-\delta))
\]
shows that \(\{|a_\delta|<r\}\) has length at most \(C\sqrt r\),
uniformly in \(\delta\). Integrating this distribution bounds the
omitted logarithmic tails uniformly by a quantity tending to zero
as \(M\to\infty\). The argument applies to every subinterval.

There is also the exact complete-prefix identity
\begin{equation}\label{lc:prefix-zero}
                   \int_0^{1/2}\log|a_\delta(t)|\,dt=0 .
\end{equation}
Shift the two cosine integrals and use the reflection and period
of \(\log|2\cos(\pi t)|\). Together they have the integral of a full
period, which is zero. Integrability includes \(\delta=0,1/2\).
We use this identity only after both actual unused products have
been bounded separately.

